\section{Conclusion}
\label{sec:conclusion}

Retrieval-state lock-in turns agreement into the wrong kind of reassurance.
When the retriever repeatedly returns the same empty or wrong state, more
samples need not test more possibilities; they may simply repeat the same
condition. The confidence question therefore changes. It is not only whether
the answer is stable, but which object is stable: the answer surface, the
retrieved evidence, or the retrieval state itself.

The experiments show that answer-state uncertainty remains a useful default,
but a bounded one. In the deployable-policy runs, $42$--$59\%$ of errors already
have zero observed answer dispersion at $N{=}5$, so an answer-only method has no
within-question disagreement signal to recover. Evidence-state diagnostics can
catch some of these failures by testing support in the retrieved text. GPS adds
a narrower graph-support view, strongest on the calibrated clinical corpus and
weak on several open-domain multi-hop suites. The field-level lesson is simple:
agreement, support, and retrieval grounding are different observables, and a
single confidence scalar hides that difference.

The practical recommendation is correspondingly conservative. Report answer-,
evidence-, and retrieval-state signals separately, and call an answer low-risk
only when all three agree: low answer uncertainty, adequate evidence support,
and present graph support. In \system{}, this conjunctive rule selects a small
subset at $91.9\%$ pooled accuracy against a $69.7\%$ base rate, and reaches
$100\%$ precision under the automated judge on the clinical target domain at
$22\%$ coverage, pending human-expert confirmation before any clinical use.
Unselected answers are not thereby wrong; they are simply not low-risk by this
audit.

The unresolved case is presence lock-in: the system retrieves a coherent but
wrong neighbourhood, and the answer is locally supported by that state.
That is the setting where relation-level path faithfulness matters most:
future systems should test not only whether an answer entity is reachable, but
whether the path uses the relation the question actually asks for. The same
diagnostic decomposition also suggests a control policy. When answer samples
agree but evidence- or retrieval-state diagnostics disagree, the retriever
should test a nearby alternative state, for example by masking the dominant
intermediate anchor while preserving the seed entities. If the answer changes,
it was anchored to that state; if it does not, the evidence was stable across
the tested perturbation.

\subsection{Limitations}
\label{sec:limitations}

The strict graph-only setting is a mechanism probe, not a deployable policy or
a one-component ablation. It changes fusion, reranking, decomposition, and
fallback together, so it demonstrates a regime in which answer-state uncertainty
fails rather than estimating production prevalence. The route logs make removed
fallback and vector augmentation the most plausible contributors in the strict
clinical collapse, but a factorial ablation would be needed to apportion the
effect.

GPS is a graph-support score in risk orientation: lower values mean stronger
support for the generated answer, not greater correctness. It is strongest on
the calibrated clinical corpus and weak on several open-domain multi-hop suites
despite the sensitivity replay in Appendix~\ref{sec:appendix_gps_sensitivity}.
Yes/no and other label-space answers remain outside its design.
The retrieval-state family is therefore justified as an observable class
(anchors, paths, subgraphs, and routes), while GPS is only one conditional scalar
instantiation of that class.

SEU is also local: it tests answer--evidence consistency, not whether the model
used the evidence. Its neutral plateau can make empty-context failures easier to
rank than populated presence-lock-in failures, where the wrong neighbourhood may
locally support the wrong answer.

The empirical scale supports the central contrasts but not fine leaderboard
claims between neighbouring metrics. Correctness labels come from an automated
semantic judge and survive an independent Llama-family re-judge, but human or
domain-expert adjudication would be needed before treating silent errors as
clinical risk. The experiments also use a single generation model
(GPT-4o-mini), and fine-grained retrieval-state diversity is logged only in the
dedicated diversity run.
The paper should therefore be read as a diagnostic methodology and fixed-system
empirical audit, not as a cross-model survey of lock-in prevalence.

The retrieval-state diagnostics are graph-native. The silent-error phenomenon
is not: dense retrieval shows it too. Extending the same decomposition to
dense-RAG retrieval states, for example through citation-path or passage-link
plausibility, is the most direct route to broader generality.

\subsection{Ethical and Safety Considerations}
\label{sec:ethics}

Nothing in these results licenses autonomous clinical answering. The
audit rule is deliberately high precision and low coverage; the unselected
majority should be routed to ordinary confidence handling, broader retrieval,
or human review. Lock-in also has an adversarial analogue: a corpus or query
can be shaped to anchor retrieval in a plausible but wrong neighbourhood.
Retrieval diversity constraints and anchor-masking controls are plausible
defences, but they remain future work.
Adversarial robustness lies outside this evaluation. The diagnostic
decomposition is still a prerequisite for it: a failure mode that cannot be
observed cannot be defended against.
